\textbf{1. Kiến thức cần nhớ}
\begin{itemize}
\item \textbf{Số trung bình (số trung bình cộng):} Cho mẫu số liệu $x_1, x_2, \ldots, x_n$.
\[ \boxed{\overline{x} = \dfrac{x_1 + x_2 + \cdots + x_n}{n}} \]
Nếu mẫu cho dưới dạng bảng tần số (giá trị $x_i$ có tần số $n_i$, cỡ mẫu $n = n_1 + \cdots + n_k$):
\[ \boxed{\overline{x} = \dfrac{n_1 x_1 + n_2 x_2 + \cdots + n_k x_k}{n}} \]
\textit{Ý nghĩa:} Đại diện cho các số liệu của mẫu, nhạy cảm với các giá trị bất thường (quá lớn hoặc quá nhỏ).
\item \textbf{Trung vị ($M_e$):} Sắp xếp mẫu số liệu theo thứ tự không giảm $x_1 \le x_2 \le \cdots \le x_n$:
\begin{itemize}
\item Nếu $n$ lẻ: $\boxed{M_e = x_{\frac{n+1}{2}}}$ (số đứng chính giữa).
\item Nếu $n$ chẵn: $\boxed{M_e = \dfrac{1}{2}\left(x_{\frac{n}{2}} + x_{\frac{n}{2}+1}\right)}$ (trung bình cộng hai số đứng giữa).
\end{itemize}
\textit{Ý nghĩa:} Dùng thay thế số trung bình khi mẫu có giá trị bất thường.
\item \textbf{Tứ phân vị ($Q_1, Q_2, Q_3$):} Chia mẫu số liệu đã sắp xếp không giảm thành $4$ phần bằng nhau (mỗi phần chứa $25\%$ số liệu):
\begin{itemize}
\item $Q_2 = M_e$ (tứ phân vị thứ hai).
\item $Q_1$ là trung vị của nửa số liệu bên trái $Q_2$ (không bao gồm $Q_2$ nếu $n$ lẻ).
\item $Q_3$ là trung vị của nửa số liệu bên phải $Q_2$ (không bao gồm $Q_2$ nếu $n$ lẻ).
\end{itemize}
\item \textbf{Mốt ($M_o$):} Là giá trị xuất hiện nhiều lần nhất (có tần số lớn nhất) trong mẫu số liệu. Một mẫu có thể có một hoặc nhiều mốt; mốt dùng được cho cả dữ liệu định tính và định lượng.
\end{itemize}

\textbf{2. Phương pháp giải các dạng}

%% DANG 0D6?3-1
\textbf{Dạng 1. Câu hỏi lý thuyết}
\begin{enumerate}
\item Đọc kĩ yêu cầu để xác định số đặc trưng được hỏi (số trung bình, trung vị, tứ phân vị hay mốt).
\item Đối chiếu định nghĩa, tính chất và ý nghĩa thực tiễn của từng số đặc trưng để chọn phát biểu đúng hoặc chọn đại lượng đại diện phù hợp.
\end{enumerate}
Lưu ý: Khi mẫu có giá trị bất thường nên dùng trung vị $M_e$ thay cho số trung bình $\overline{x}$; với dữ liệu định tính hoặc cỡ giày, cỡ áo nên dùng mốt $M_o$.

%% DANG 0D6?3-2
\textbf{Dạng 2. Số trung bình cộng}
\begin{enumerate}
\item Xác định cỡ mẫu $n$ (đếm số phần tử hoặc tính tổng các tần số $n_1 + n_2 + \cdots + n_k$).
\item Áp dụng công thức $\overline{x} = \dfrac{\sum x_i}{n}$ hoặc $\overline{x} = \dfrac{\sum n_i x_i}{n}$ để tính toán (có thể dùng chế độ thống kê trên máy tính cầm tay).
\end{enumerate}
Lưu ý: Khi mẫu cho dưới dạng bảng tần số, tránh nhầm lẫn chia cho số giá trị khác nhau $k$ thay vì chia cho cỡ mẫu $n$.

%% DANG 0D6?3-3
\textbf{Dạng 3. Số trung vị}
\begin{enumerate}
\item Sắp xếp tất cả các giá trị của mẫu số liệu theo thứ tự không giảm (giữ nguyên các giá trị lặp lại).
\item Xác định cỡ mẫu $n$ là chẵn hay lẻ, từ đó tìm vị trí chính giữa và tính $M_e$ theo công thức.
\end{enumerate}
Lưu ý: Sai lầm phổ biến nhất là quên sắp xếp mẫu số liệu theo thứ tự không giảm trước khi lấy số ở giữa.

%% DANG 0D6?3-4
\textbf{Dạng 4. Tứ phân vị}
\begin{enumerate}
\item Sắp xếp mẫu số liệu theo thứ tự không giảm và tìm trung vị $Q_2 = M_e$.
\item Tách mẫu thành hai nửa: nửa bên trái $Q_2$ và nửa bên phải $Q_2$ (nếu $n$ lẻ thì bỏ qua phần tử chính giữa $Q_2$).
\item Tìm trung vị của nửa bên trái được $Q_1$ và trung vị của nửa bên phải được $Q_3$.
\end{enumerate}
Lưu ý: Khi $n$ lẻ, tuyệt đối không đưa giá trị $Q_2$ vào hai nửa số liệu khi tính $Q_1$ và $Q_3$; luôn kiểm tra điều kiện $Q_1 \le Q_2 \le Q_3$.

%% DANG 0D6?3-5
\textbf{Dạng 5. Mốt}
\begin{enumerate}
\item Lập bảng tần số (hoặc đếm số lần xuất hiện của từng giá trị khác nhau trong mẫu).
\item Tìm tần số lớn nhất và kết luận các giá trị tương ứng là mốt $M_o$.
\end{enumerate}
Lưu ý: Mốt là \textit{giá trị} của mẫu số liệu, không phải là \textit{tần số} lớn nhất; một mẫu số liệu có thể có nhiều hơn một mốt.
